\begin{frame}{정보이득: 나누기 전 엔트로피 $-$ 가중평균}
  \kb{엔트로피}(entropy, Week 2.3 섀넌 정보이론):
  \[ H = -\sum_{k=1}^{K} p_k \log_2 p_k \]
  지니불순도와 마찬가지로 순수할수록 0, 섞일수록 크다.
  \vskip0.6em
  노드를 왼쪽($L$)과 오른쪽($R$)으로 나눴을 때 \kb{정보이득}(Information Gain):
  \[ \text{IG} = H(\text{parent})
    - \left( \tfrac{|L|}{|L|+|R|} H(L) + \tfrac{|R|}{|L|+|R|} H(R) \right) \]
  \begin{itemize}
    \item IG가 클수록 그 질문이 데이터를 잘 나눴다는 뜻
    \item 트리를 만들 때 각 단계에서 \textbf{IG(또는 지니 감소량)가 가장 큰
      질문}을 고른다 — 스무고개의 ``가장 잘 좁히는 질문을 먼저''이
      알고리즘이 된 것
  \end{itemize}
\end{frame}
